\noindent\textbf{\textit{Résumé}---}\textbf{\textit{Les affirmations sur la rapidité des langages de programmation
reposent rarement sur des mesures d'applications complètes. Nous comparons cinq implémentations équivalentes d'une
même API REST de gestion de tâches adossée à SQLite (Rust, Go, PHP, Node.js et Python) sous charge croissante (10, 50
et 100 utilisateurs simultanés), sur cinq fonctionnalités (création, liste paginée, recherche, jointure, statistiques).
Chaque cas est répété 30 fois, soit 2\,250 essais et 1,8 million de requêtes ; nous mesurons le débit, les latences
(médiane, p95, p99), le temps CPU par requête, la mémoire et le temps de démarrage. Sur les fonctionnalités légères,
Rust traite jusqu'à 9,3 fois plus de requêtes par seconde que Python (13\,373 contre 1\,433 en création) ; sur la
recherche et les statistiques, où SQLite domine, l'écart tombe à 1,5 et à environ 1. Ces résultats valent pour ces
implémentations, sur un seul cœur ; ils ne se généralisent pas aux langages.}}

\smallskip
\noindent\textit{\textbf{Mots-clés}---}performance, test de charge, API REST, SQLite, latence, Rust, Go, PHP, Node.js,
Python
